Kurs:Algebraische Kurven (Osnabrück 2025-2026)/Arbeitsblatt 12/latex
Kurs:Algebraische Kurven (Osnabrück 2025-2026)/Arbeitsblatt 12/latex

\setcounter{section}{12}


  \zwischenueberschrift{Übungsaufgaben}


\inputaufgabe

{}

{


Man erläutere das Konzept eines Punktes in der euklidischen
\zusatzklammer {koordinatenfreien} {} {}
Geometrie und in der kartesischen Geometrie. In welchen Situationen ist der Übergang zu Koordinaten sinnvoll?


}

{} {}


\inputaufgabe

{}

{


Was ist ein Punkt in der algebraischen Geometrie? Welche Punktkonzepte haben Sie in dem Kurs über algebraische Kurven kennengelernt und wie hängen diese miteinander zusammen?


}

{} {}


\inputaufgabe

{}

{


Bestimme das
$K$-\definitionsverweis {Spektrum}{}{}
von

\mathl{K^d}{.}


}

{} {}


\inputaufgabe

{}

{


Bestimme das
$\R$-\definitionsverweis {Spektrum}{}{}
der
$\R$-\definitionsverweis {Algebra}{}{}
${\mathbb C}$.


}

{} {}


\inputaufgabe

{}

{


Bestimme das
$\R$-\definitionsverweis {Spektrum}{}{}
von

\mathl{\R[X,Y]/(X^2+Y^2-1)}{.}


}

{} {}


\inputaufgabegibtloesung

{}

{


Es sei $K$ ein
\definitionsverweis {algebraisch abgeschlossener Körper}{}{}
und $R$ eine kommutative
$K$-\definitionsverweis {Algebra von endlichem Typ}{}{.}
Zeige, dass die Punkte aus

\mathl{K\!-\!\operatorname{Spek}\, \left(  R  \right)}{} den
\definitionsverweis {maximalen Idealen}{}{}
in $R$ entsprechen.


}

{} {}


\inputaufgabe

{}

{


Es sei $R$ eine kommutative $K$-Algebra von endlichem Typ. Zeige, dass für jedes
\definitionsverweis {Ideal}{}{}


\mavergleichskette

{\vergleichskette

{ {\mathfrak a}
}

{ \subseteq }{ R
}

{  }{
}

{    }{
}

{   }{
}

}
{}{}{}
in

\mathl{K\!-\!\operatorname{Spek}\, \left(   R   \right)}{} die Gleichheit


\mavergleichskettedisp

{\vergleichskette

{ V ( {\mathfrak a} )
}

{ =} { V(\operatorname{rad} \,  ( {\mathfrak a} ))
}

{ } {
}

{ } {
}

{ } {
}

}
{}{}{}
gilt.


}

{} {}


\inputaufgabe

{}

{


Zeige, dass die Zariski-Topologie auf dem $K$-Spektrum einer endlich erzeugten kommutativen $K$-Algebra $R$ wirklich eine \definitionsverweis {Topologie}{}{}
ist.


}

{} {}


\inputaufgabe

{}

{


Es sei $K$ ein
\definitionsverweis {Körper}{}{} und sei


\mavergleichskette

{\vergleichskette

{V
}

{ \subseteq }{ { {\mathbb A}_{  K }^{  n   } }
}

{  }{
}

{    }{
}

{   }{
}

}
{}{}{}
eine
\definitionsverweis {affin-algebraische Menge}{}{}
mit
\definitionsverweis {Verschwindungsideal}{}{}


\mathl{\operatorname{Id} \, \left(   V   \right)}{} und
\definitionsverweis {Koordinatenring}{}{}


\mavergleichskettedisp

{\vergleichskette

{R
}

{ \defeq} { K[X_1 , \ldots , X_n]/ \operatorname{Id} \, \left(   V   \right)
}

{ } {
}

{ } {
}

{ } {
}

}
{}{}{.}
Zeige, dass das
$K$-\definitionsverweis {Spektrum}{}{}
von $R$
\definitionsverweis {homöomorph}{}{}
zu $V$ ist.


}

{} {}


Sämtliche Versionen des Hilbertschen Nullstellensatz wie
Korollar 11.11
übertragen sich auf das
$K$-\definitionsverweis {Spektrum}{}{}
einer endlich erzeugten $K$-Algebra über einem algebraisch abgeschlossenen Körper $K$.


\inputaufgabe

{}

{


Es sei $K$ ein
\definitionsverweis {algebraisch abgeschlossener Körper}{}{}
und sei $R$ eine
$K$-\definitionsverweis {Algebra von endlichem Typ}{}{.}
Beweise, dass zwischen den
\definitionsverweis {abgeschlossenen Teilmengen}{}{}
des
$K$-\definitionsverweis {Spektrums}{}{}


\mathl{K\!-\!\operatorname{Spek}\, \left(   R   \right)}{} und den
\definitionsverweis {Radikalen}{}{}
in $R$ eine bijektive Korrespondenz besteht.


}

{} {}


\inputaufgabe

{}

{


Es sei $K$ ein
\definitionsverweis {algebraisch abgeschlossener Körper}{}{}
und sei $R$ eine
\definitionsverweis {reduzierte}{}{}
$K$-\definitionsverweis {Algebra von endlichem Typ}{}{.}
Beweise den \stichwort {Identitätssatz} {} in der folgenden Gestalt: Wenn für


\mavergleichskette

{\vergleichskette

{ f,g
}

{ \in }{ R
}

{  }{
}

{    }{
}

{   }{
}

}
{}{}{}
gilt, dass


\mavergleichskette

{\vergleichskette

{ f(P)
}

{ = }{ g(P)
}

{  }{
}

{    }{
}

{   }{}

}
{}{}{}
ist für alle


\mavergleichskette

{\vergleichskette

{ P
}

{ \in }{ K\!-\!\operatorname{Spek}\, \left(  R  \right)
}

{  }{
}

{    }{
}

{   }{
}

}
{}{}{,}
so ist


\mavergleichskette

{\vergleichskette

{ f
}

{ = }{ g
}

{  }{
}

{    }{
}

{   }{
}

}
{}{}{.}


}

{} {}


\inputaufgabegibtloesung

{}

{


Es sei $K$ ein
\definitionsverweis {Körper}{}{,}
$R$ eine
\definitionsverweis {endlich erzeugte}{}{}
$K$-Algebra, sei


\mavergleichskette

{\vergleichskette

{ {\mathfrak a}
}

{ \subseteq }{ R
}

{  }{
}

{    }{
}

{   }{
}

}
{}{}{}
ein
\definitionsverweis {Ideal}{}{}
und sei


\mavergleichskette

{\vergleichskette

{ X
}

{ = }{ K\!-\!\operatorname{Spek}\, \left(  R  \right)
}

{  }{
}

{    }{
}

{   }{
}

}
{}{}{.}
In welcher Beziehung stehen die beiden Aussagen


\mathdisp {V( {\mathfrak a}) = \emptyset  \text{ und } {\mathfrak a} \text{ ist das Einheitsideal}} {  }


und die beiden Aussagen


\mathdisp {V( {\mathfrak a}) = X  \text{ und } {\mathfrak a} \text{ ist nilpotent}} {  }


zueinander. Zeige, dass die Antwort davon abhängt, ob $K$
\definitionsverweis {algebraisch abgeschlossen}{}{}
ist oder nicht.


}

{} {}


\inputaufgabe

{}

{


Man beschreibe zu einer kommutativen
$K$-\definitionsverweis {Algebra}{}{}
$R$ von endlichem Typ die
\definitionsverweis {Spektrumsabbildung}{}{,}
die zum
\definitionsverweis {Strukturhomomorphismus}{}{}
\maabb {} {K} {R
} {}
der Algebra gehört.


}

{} {}


\inputaufgabe

{}

{


Es seien

\mathl{R,S,T}{} kommutative
$K$-\definitionsverweis {Algebren von endlichem Typ}{}{}
und
\maabb {\varphi} { R } { S
} {}
und
\maabb {\psi} { S } { T
} {}
seien
$K$-\definitionsverweis {Algebrahomomorphismen}{}{.}
Man zeige, dass für die zugehörigen
\definitionsverweis {Spektrumsabbildungen}{}{}


\mavergleichskettedisp

{\vergleichskette

{ ( \psi \circ  \varphi )^*
}

{ =} { \varphi^*  \circ  \psi^*
}

{ } {
}

{ } {
}

{ } {
}

}
{}{}{}
gilt. Ferner zeige man, dass zur Identität
\maabb {\operatorname{Id}} { R } { R
} {}
auch $\operatorname{Id}^*$ die Identität ist.


}

{} {}


\inputaufgabe

{}

{


Man gebe ein Beispiel von zwei kommutativen
$K$-\definitionsverweis {Algebren}{}{}
$R,S$
\definitionsverweis {von endlichem Typ}{}{}
und einer
\definitionsverweis {stetigen Abbildung}{}{}
zwischen den zugehörigen
$K$-\definitionsverweis {Spektren}{}{,}
die nicht von einem
$K$-\definitionsverweis {Algebrahomomorphismus}{}{}
herrühren kann.


}

{} {}


\inputaufgabe

{}

{


Es sei $K$ ein
\definitionsverweis {Körper}{}{}
und $R$ eine kommutative
$K$-\definitionsverweis {Algebra von endlichem Typ}{}{,}
und sei


\mavergleichskette

{\vergleichskette

{ F
}

{ \in }{ R
}

{  }{
}

{    }{
}

{   }{
}

}
{}{}{.}
Es sei
\maabbdisp {\varphi^*} { K\!-\!\operatorname{Spek}\, \left(   R   \right) } { {\mathbb A}^{1}_{ K }
} {}
die zum
\definitionsverweis {Einsetzungshomomorphismus}{}{}
gehörende
\definitionsverweis {Spektrumsabbildung}{}{.}
Zeige, dass


\mavergleichskettedisp

{\vergleichskette

{ (\varphi^*)^{-1} (0)
}

{ =} { V(F)
}

{ } {
}

{ } {
}

{ } {
}

}
{}{}{}
ist.


}

{} {}


\inputaufgabe

{}

{


Es sei $K$ ein
\definitionsverweis {Körper}{}{}
und sei $R$ eine kommutative
$K$-\definitionsverweis {Algebra von endlichem Typ}{}{}
mit der
\definitionsverweis {Reduktion}{}{}


\mavergleichskette

{\vergleichskette

{ S
}

{ = }{ R_{\rm red}
}

{  }{
}

{    }{
}

{   }{
}

}
{}{}{.}
Zeige, dass es eine natürliche
\definitionsverweis {Homöomorphie}{}{}


\mavergleichskettedisp

{\vergleichskette

{ K\!-\!\operatorname{Spek}\, \left(   R   \right)
}

{ \cong} { K\!-\!\operatorname{Spek}\, \left(   S   \right)
}

{ } {
}

{ } {
}

{ } {
}

}
{}{}{}
gibt.


}

{} {}


\inputaufgabe

{}

{


Es sei $K$ ein
\definitionsverweis {Körper}{}{}
und seien


\mavergleichskette

{\vergleichskette

{ F_i
}

{ \in }{ K[X_1 , \ldots , X_n]
}

{  }{
}

{    }{
}

{   }{
}

}
{}{}{}
\definitionsverweis {Polynome}{}{}
für


\mavergleichskette

{\vergleichskette

{ i
}

{ = }{ 1 , \ldots , m
}

{  }{
}

{    }{
}

{   }{
}

}
{}{}{.}
Es sei
\maabbeledisp {\varphi} { K[Y_1 , \ldots , Y_m] } { K[X_1 , \ldots , X_n]
} { Y_i } { F_i
} {,}
der zugehörige
\definitionsverweis {Einsetzungshomomorphismus}{}{.}
Zeige, dass die
\definitionsverweis {Spektrumsabbildung}{}{}
\maabbdisp {\varphi^*} { { {\mathbb A}_{  K  }^{  n   } } = K\!-\!\operatorname{Spek}\, \left(   K[X_1 , \ldots , X_n]   \right)  } { { {\mathbb A}_{  K }^{  m   } }=K\!-\!\operatorname{Spek}\, \left(   K[Y_1 , \ldots , Y_m]   \right)
} {}
\zusatzklammer {über die Identifizierung aus
Lemma 12.3} {} {}
mit der direkten
\definitionsverweis {polynomialen Abbildung}{}{}


\mathdisp {\left(  x_1 , \ldots , x_n  \right)  \longmapsto \left(  F_1 \left(  x_1 , \ldots , x_n  \right)  , \ldots , F_m \left(  x_1 , \ldots , x_n  \right)   \right)} {  }


übereinstimmt.


}

{} {}


\inputaufgabe

{}

{


Welche \anfuehrung{Funktoren}{} in der Mathematik kennen Sie?


}

{} {}


Bei den folgenden Aufgaben betrachten wir zu einem beliebigen topologischen Raum
\zusatzklammer {Mannigfaltigkeit, Teilmenge des $\R^n$, reelles Intervall} {} {}
den Ring der stetigen reellwertigen Funktionen darauf. Dabei sollte man die Räume in Analogie zu den $K$-Spektren und die Funktionenringe in Analogie zu den
\zusatzklammer {Koordinaten} {} {-}
ringen sehen.


\inputaufgabe

{}

{


Es sei $X$ ein
\definitionsverweis {topologischer Raum}{}{}
und


\mavergleichskettedisp

{\vergleichskette

{R
}

{ =} { C(X, \R)
}

{ =} { \left\{ f:X \longrightarrow \R  \mid  f \text{ stetige Abbildung}         \right\}
}

{ } {
}

{ } {
}

}
{}{}{.}
Zeige, dass $R$ ein
\definitionsverweis {kommutativer Ring}{}{}
ist.


}

{} {}


\inputaufgabe

{}

{


Wir betrachten den Ring


\mavergleichskette

{\vergleichskette

{R
}

{ = }{ C(\R,\R)
}

{  }{
}

{    }{
}

{   }{
}

}
{}{}{}
der
\definitionsverweis {stetigen Funktionen}{}{}
von $\R$ nach $\R$. Handelt es sich um einen
\definitionsverweis {Integritätsbereich}{}{?}


}

{} {}


\inputaufgabe

{}

{


Es sei


\mavergleichskette

{\vergleichskette

{T
}

{ \subseteq }{ \R
}

{  }{
}

{    }{
}

{   }{
}

}
{}{}{}
eine Teilmenge. Zeige, dass im Ring der
\definitionsverweis {stetigen Funktionen}{}{}


\mavergleichskettedisp

{\vergleichskette

{R
}

{ =} { C(\R,\R)
}

{ } {
}

{ } {
}

{ } {
}

}
{}{}{}
die Teilmenge


\mavergleichskettedisp

{\vergleichskette

{I
}

{ =} { \left\{ f \in R  \mid   f(x) = 0 \text{ für alle } x \in T         \right\}
}

{ } {
}

{ } {
}

{ } {
}

}
{}{}{}
ein
\definitionsverweis {Ideal}{}{}
in $R$ ist.


}

{} {}


\inputaufgabe

{}

{


Wir betrachten das Ideal zu


\mavergleichskette

{\vergleichskette

{T
}

{ = }{ \{0\}
}

{ \subseteq }{\R
}

{    }{
}

{   }{
}

}
{}{}{}
im Sinne von
Aufgabe 12.22.
Ist dies ein
\definitionsverweis {Hauptideal}{}{?}


}

{} {}


\inputaufgabe

{}

{


Es sei $X$ ein
\definitionsverweis {topologischer Raum}{}{}
und


\mavergleichskette

{\vergleichskette

{ R
}

{ = }{ \operatorname{C}^0 \, ( X , \R)
}

{  }{
}

{    }{
}

{   }{
}

}
{}{}{}
der
\definitionsverweis {Ring der stetigen Funktionen}{}{}
auf $X$. Es sei


\mavergleichskette

{\vergleichskette

{ T
}

{ \subseteq }{ X
}

{  }{
}

{    }{
}

{   }{
}

}
{}{}{}
eine Teilmenge. Zeige, dass die Teilmenge


\mavergleichskettedisp

{\vergleichskette

{ I
}

{ =} { \left\{ f \in R  \mid  f\vert_T=0        \right\}
}

{ } {
}

{ } {
}

{ } {
}

}
{}{}{}
ein
\definitionsverweis {Ideal}{}{}
in $R$ ist. Definiere einen
\definitionsverweis {Ringhomomorphismus}{}{}
\maabbdisp {} { R/I } { \operatorname{C}^0 \, ( T , \R)
} {.}
Ist dieser immer injektiv? Surjektiv?


}

{} {}


\inputaufgabe

{}

{


Es seien
\mathkor {} {X} {und} {Y} {}
\definitionsverweis {topologische Räume}{}{} und
\maabbdisp {\varphi} { X } { Y
} {}
eine
\definitionsverweis {stetige Abbildung}{}{.}
Zeige, dass dies einen
\definitionsverweis {Ringhomomorphismus}{}{}
\maabbeledisp {} { C(Y, \R)} {C(X, \R)
} { f } { f\circ \varphi
} {,}
induziert.


}

{} {}


\inputaufgabe

{}

{


Es sei $X$ eine Teilmenge von $\R$ und $C(X, \R)$ der
\definitionsverweis {Ring der stetigen Funktionen}{}{}
von $X$ nach $\R$.
Dann ist durch
\maabbeledisp {\varphi} { C(\R, \R)  } { C(X, \R)
} { f } { f \vert_X
} {,}
ein
\definitionsverweis {Ringhomomorphismus}{}{}
gegeben.
\aufzaehlungzwei {Zeige, dass $\varphi$ genau dann
\definitionsverweis {surjektiv}{}{}
ist, wenn $X$
\definitionsverweis {abgeschlossen}{}{}
ist.
} {Für welche Mengen $X$ ist $\varphi$
\definitionsverweis {injektiv}{}{?}
}


}

{} {}


  \zwischenueberschrift{Aufgaben zum Abgeben}


\inputaufgabe

{4 (2+2)}

{


Es sei $K$ ein unendlicher
\definitionsverweis {Körper}{}{}
und $R$ eine kommutative
$K$-\definitionsverweis {Algebra von endlichem Typ}{}{,}
und sei


\mavergleichskette

{\vergleichskette

{ F
}

{ \in }{ R
}

{  }{
}

{    }{
}

{   }{
}

}
{}{}{.}
Es sei
\maabbdisp {\varphi^*} { K\!-\!\operatorname{Spek}\, \left(   R   \right) } { {\mathbb A}^{1}_{ K }
} {}
die zum
\definitionsverweis {Einsetzungshomomorphismus}{}{}
gehörende
\definitionsverweis {Spektrumsabbildung}{}{.}
\aufzaehlungzweiabc{Zeige, dass $F$ genau dann konstant ist, wenn $\varphi^*$ konstant ist.
}{Zeige, dass diese Aussage für einen endlichen Körper nicht gelten muss.
}


}

{Man mache sich dabei auch die unterschiedlichen Bedeutungen von \anfuehrung{konstant}{} klar.} {}


\inputaufgabe

{4}

{


Man gebe ein Beispiel von zwei
\definitionsverweis {integren}{}{}
${\mathbb C}$-\definitionsverweis {Algebren von endlichem Typ}{}{}
$R$ und $S$ und einem
${\mathbb C}$-\definitionsverweis {Algebrahomomorphismus}{}{}
\maabb {\varphi} { R } { S
} {,}
der kein
\definitionsverweis {Ringisomorphismus}{}{}
ist, wo aber die induzierte
\definitionsverweis {Spektrumsabbildung}{}{}
\maabb {\varphi^*} { {\mathbb C}\!-\!\operatorname{Spek}\, \left(   S   \right) } { {\mathbb C}\!-\!\operatorname{Spek}\, \left(   R   \right)
} {}
ein
\definitionsverweis {Homöomorphismus}{}{}
ist.


}

{} {}


\inputaufgabe

{3}

{


Es sei $K$ ein
\definitionsverweis {algebraisch abgeschlossener Körper}{}{}
und sei $R$ eine $K$-Algebra von endlichem Typ. Wir betrachten eine
\definitionsverweis {endliche}{}{}
Erweiterung der Form
\maabbdisp {\varphi} { R } { S = R [X]/ \left(  X^n + r_{n-1} X^{n-1} + \cdots + r_2X^2 +r_1X +r_0  \right)
} {.}
Zeige, dass dann
\maabb {\varphi^*} { K\!-\!\operatorname{Spek}\, \left(  S  \right) } { K\!-\!\operatorname{Spek}\, \left(  R  \right)
} {} surjektiv ist.


}

{} {}


\inputaufgabe

{5}

{


Betrachte das
\definitionsverweis {Ideal}{}{}


\mavergleichskettedisp

{\vergleichskette

{ {\mathfrak a}
}

{ =} { \left(  U^5-V^3, U^{11}-W^3,V^{11}-W^5  \right)
}

{ \subseteq} { K[U,V,W]
}

{ } {
}

{ } {
}

}
{}{}{}
und das zugehörige
\definitionsverweis {Nullstellengebilde}{}{}


\mavergleichskette

{\vergleichskette

{ Z
}

{ = }{ V({\mathfrak a} )
}

{ \subseteq }{ { {\mathbb A}_{ K }^{  3   } }
}

{    }{
}

{   }{
}

}
{}{}{.}
Zeige, dass

\mathl{W-U^2V}{} zum Radikal von ${\mathfrak a}$ gehört. Zeige damit, dass $Z$ isomorph zu einer ebenen algebraischen Kurve ist.


}

{} {Man benutze, dass das Radikal der Durchschnitt der Primideale ist, die es umfassen, oder man reduziere auf den Fall, dass $K$ algebraisch abgeschlossen ist.}
